% Generated from locale/tr/content/proof-theory/sequent-calculus/translations.tex; do not hand-edit.


\olfileid[tr]{pt}{seq}{tr}

\olsection{\Log{G1c} ile \Log{G3c} Arasında Çeviri}

Hem \Log{G1c}'nin hem de \Log{G3c}'nin klasik mantık için sağlam ve tam
olduğunu, yani her birinin bütün !!{structure}s{}da doğru olan sekantları ve
yalnızca onları ispatladığını belirtmiştik. Özellikle aynı sekantları
ispatlarlar. Artık bu olguyu sağlamlık ve tamlık teoremleri üzerinden
dolanmadan, salt ispat kuramsal yöntemlerle gösterebiliriz. Böyle bir ispat,
\Log{G1c}'deki !!{proof}s{}ı \Log{G3c}'deki !!{proof}s{}a ve ters yönde
dönüştürme yolları sağlar.

\begin{prop}
$\Log{G1c} \Proves S$ ise $\Log{G3c} \Proves S$.
\end{prop}

\begin{proof}
  $\pi$, \Log{G1c}'de !!a{proof} ise \Log{G3c}'de $S$'nin bir
  !!a{proof}{}ı~$\pi'$ bulunduğunu gösteriyoruz. Bunu $n = \pheight{\pi}$
  üzerinde tümevarımla yaparız.

  $n = 0$ ise $\pi$ bir aksiyomdur. $\lfalse \Sequent \quad$ aksiyomu
  \Log{G3c}'de de bir aksiyomdur (bağlamlar $\Gamma, \Delta$ boş alınır).
  \Log{G3c}'de gelişigüzel $!A$ için $!A \Sequent !A$ aksiyomlarına izin
  verilmez; aksiyomlarda $!A$ atomik olmalıdır. Ancak \olref[ex]{prop:G3c-axioms}
  içinde bütün bu sekantların \Log{G3c}'de !!{provable} olduğunu
  göstermiştik.
  
  $n>0$ ise $\pi$ bir~$R$ kuralıyla biter. Tümevarım varsayımı, bu kuralın
  öncüllerinin \Log{G3c}'de !!{proof}s{}ı bulunduğunu, yani doğrudan
  alt-!!{proof}s{}ın \Log{G3c}'ye çevirileri olduğunu güvence altına alır.
  $R$ kuralı \Log{G3c}'nin de bir kuralıysa onu öncüllerin çevrilmiş
  !!{proof}s{}ına uygularız. \Log{G3c}'nin kuralı olmayan kuralların en azından
  kabul edilebilir olduğunu daha önce göstermiştik. Dolayısıyla $R$ kuralı
  $\LeftR{\land}$ ya da $\RightR{\lor}$ ise
  \olref[adm]{prop:lor-G3-adm}; zayıflatma kuralıysa
  \olref[adm]{prop:weak-G3c-adm}; büzülme kuralıysa
  \olref[inv]{prop:G3c-cont-adm} uygulanır.
\end{proof}

\begin{prop}
$\Log{G3c} \Proves S$ ise $\Log{G1c} \Proves S$.
\end{prop}

\begin{proof}
  Yine $S$'nin bir !!a{proof}{}ı olan~$\pi$'nin !!{height} değeri üzerinde tümevarımla
  ilerleriz. \Log{G3c}'nin her aksiyomu, \Log{G1c}'de bir aksiyomdan
  başlayıp \LeftR{\Weakening} ve \RightR{\Weakening} kurallarını birkaç kez
  uygulayarak !!{prove}{}nabilir:
  \[
  \Axiom$!C \fCenter !C$
  \RightLabel{\LeftR{\Weakening}}
  \Deduce$!C, \Gamma \fCenter !C$
  \RightLabel{\RightR{\Weakening}}
  \Deduce$!C, \Gamma \fCenter \Delta, !C$
  \DisplayProof
\qquad
  \Axiom$\lfalse \fCenter$
  \RightLabel{\LeftR{\Weakening}}
  \Deduce$\lfalse, \Gamma \fCenter$
  \RightLabel{\RightR{\Weakening}}
  \Deduce$\lfalse, \Gamma \fCenter \Delta$
  \DisplayProof
  \]
  $\pi$ bir~$R$ kuralıyla bitiyorsa tümevarım varsayımı öncüllerin
  \Log{G1c}'de !!{proof}s{}ını verir ve aynı~$R$ kuralını bu !!{proof}s{}a
  uygulayabiliriz. İstisnalar, \Log{G1c} ile \Log{G3c}'de farklı olan
  \LeftR{\land} ve \RightR{\lor} kurallarıdır. \Log{G3c} sürümleri
  \Log{G1c}'de şöyle türetilebilir:
  \[
  \Axiom$!A, !B, \Gamma \fCenter \Delta$
  \RightLabel{\LeftR{\land}}
  \UnaryInf$!A \land !B, !B, \Gamma \fCenter \Delta$
  \RightLabel{\LeftR{\land}}
  \UnaryInf$!A \land !B, !A \land !B, \Gamma \fCenter \Delta$
  \RightLabel{\LeftR{\Contraction}}
  \UnaryInf$!A \land !B, \Gamma \fCenter \Delta$
    \DisplayProof
\qquad
  \Axiom$\Gamma \fCenter \Delta, !A, !B$
  \RightLabel{\RightR{\lor}}
  \UnaryInf$\Gamma \fCenter \Delta, !A \lor !B, !B$
  \RightLabel{\RightR{\lor}}
  \UnaryInf$\Gamma \fCenter \Delta, !A \lor !B, !A \lor !B$
  \RightLabel{\RightR{\Contraction}}
  \UnaryInf$\Gamma \fCenter \Delta, !A \lor !B$
    \DisplayProof
  \]
\end{proof}

